\section{Observation Contract}

An observation request is prospective: it states which evidence and derived
capabilities a caller needs before execution. Preparation resolves this
request against the bundle's catalog, checks required prerequisites, and
records any permitted optional downgrade. An \texttt{EvidenceReceipt} is
retrospective: it records which requested channels and capabilities were
fulfilled, incomplete, or failed after the run. These objects answer
different questions and are intentionally separate.

The emission model preserves channel distinctions:

\begin{description}[leftmargin=!,labelwidth=0.25\linewidth]
  \item[Normalized events] Machine-independent temporal and control facts
        selected for the run.
  \item[Native evidence] Typed machine-specific records whose schema and
        interpretation remain bundle-owned.
  \item[Frame artifacts] Requested raw frame data or a requested fingerprint
        of frame content.
  \item[Snapshots] Machine-state artifacts exposed at negotiated snapshot
        points; a snapshot is not itself a replay package.
  \item[Input-value evidence] Values of selected scheduled or live inputs,
        with their declared application coordinates.
  \item[Capability outputs] Derived interpretations produced by a bundle's
        normalizer from declared evidence dependencies.
  \item[Capability receipts] Per-capability accounting for availability and
        fulfillment.
  \item[Run lifecycle] Events that delimit preparation, execution, and run
        completion or termination.
  \item[Evidence receipt] Aggregate retrospective accounting for requested
        evidence and capabilities.
\end{description}

The distinction between event and artifact is important. \texttt{FrameCompleted}
is a temporal and control fact that a frame boundary was reached.
\texttt{FrameArtifact} is raw or fingerprinted frame evidence requested by
the caller. A capability output is a derived interpretation of evidence, not
the primary evidence from which it was derived. Receipts account for
fulfillment and persistence; they are not substitutes for evidence records.

Normalizers operate as streaming reducers: \texttt{begin}, zero or more calls
to \texttt{observe(emission)}, and \texttt{finalize} with the optional
negotiated final state. Each reducer retains the state required by its declared
capability and selected inputs. GlassVM does not silently retain all emissions
on behalf of a reducer, and does not construct a canonical whole-run trace.
The selected channels and budgets define what is delivered; the consumer
decides whether to retain it beyond the synchronous callback.

Evidence budgets, where configured, constrain selected channels during
production. Their fulfillment is visible in the receipt and remains separate
from machine termination. Persistence has a second accounting boundary in
the recorder receipt: evidence emitted and evidence stored are related but
not identical quantities.
